%% Uji opsi mode=proposal: memastikan spasi berubah ke 1,15 dan label jenis TA berbeda.
%% Kompilasi: pdflatex -interaction=nonstopmode uji-proposal.tex  (2x)
\documentclass[jenis=prototipe,mode=proposal,jenjang=diploma,sitasi=ieee,logo=logo-placeholder]{unnes-ta}

\identitas{
  judul    = {PENGUJIAN MODE PROPOSAL},
  nama     = {Nama Mahasiswa},
  nim      = {1234567890},
  prodi    = {Teknik Elektronika},
  fakultas = {Fakultas Teknik},
  gelar    = {Diploma Tiga},
  tahun    = {2026},
}

\begin{document}
\bagianawal
\chapter*{UJI MODE PROPOSAL}
Halaman ini memakai mode proposal sehingga jarak barisnya 1,15 spasi.

\bagianutama
\chapter{PENDAHULUAN}
\section{Latarbelakang}
Uji spasi proposal. Baris ini panjang dengan sengaja supaya jarak antarbaris dapat diukur
secara terprogram dari koordinat teks pada berkas PDF hasil kompilasi, sehingga perbedaan
1,5 spasi dan 1,15 spasi dapat dibuktikan secara objektif, bukan hanya dilihat mata.
Kata \textit{longtable} dan table istilah asing diuji agar pemenggalan kata terlihat.
\end{document}
